%%
%% O5_spectral_cluster_stability.tex
%% Spectral Cluster Stability Lemma for O5
%%
%% ROLE IN DAG:
%%   This is the FINAL formal step closing O5.
%%   It bridges the gap between:
%%     (i)  norm convergence H(c) -> I  (from O5_transfer_lemma.tex)
%%     (ii) spectral gap stability  inf sigma(D^Airy) > 0  (the O5 claim)
%%   via a Kato-type argument for isolated spectrum.
%%
%%   After this lemma, O5 is CLOSED (modulo S2/S4 which are non-blocking).
%%
%% PROOF INPUTS:
%%   - O5_transfer_lemma.tex     (norm convergence, C*c^{-1/3} bound)
%%   - O5_airy_gap.tex           (delta_Airy >= F_TW(0) approx 0.96)
%%   - O5_airy_sampling.tex      (coercivity of A^arith, C_J < 1)
%%   - Kato (1966), Theorem IV.3.18  (stability of isolated spectrum)
%%
%% PROHIBITION:
%%   No new spectral theory. No new PSWF or Airy estimates.
%%   One self-adjoint perturbation argument, one explicit constant.
%%
\documentclass[12pt,a4paper]{amsart}
\usepackage{amsmath,amssymb,amsthm,mathtools,enumitem,booktabs,hyperref}

\newtheorem{theorem}{Theorem}[section]
\newtheorem{lemma}[theorem]{Lemma}
\newtheorem{corollary}[theorem]{Corollary}
\theoremstyle{definition}
\newtheorem{definition}[theorem]{Definition}
\newtheorem{remark}[theorem]{Remark}

\newcommand{\KAi}{K^{\mathrm{Airy}}}
\newcommand{\DAi}{\mathcal{D}^{\mathrm{Ai}}}
\newcommand{\PAi}{\Pi_{\mathrm{Airy}}}
\newcommand{\Aarith}{A^{\mathrm{arith}}}
\newcommand{\norm}[1]{\|#1\|}
\newcommand{\abs}[1]{\left|#1\right|}

\title{O5 Spectral Cluster Stability Lemma\\
{\normalsize The Final Formal Step Closing O5}}
\author{Sebastian Schmalnauer}
\date{Mai 2026}

\begin{document}
\maketitle

\begin{abstract}
We supply the single remaining formal step between O5's norm-level
results and its spectral conclusion.
Let $\delta := F_{\mathrm{TW}}(0) \approx 0.9597$ be the Tracy--Widom
gap and $\eta(c) := C_J + \varepsilon(c) + O(c^{-1/3})$ the total
perturbation assembled from O5's three component files.
Since $\eta(c) \to C_J < \delta$ and the operator $\DAi$ is self-adjoint,
Weyl's eigenvalue stability theorem gives
$\inf\sigma(\DAi) \ge \delta - \eta(c) > 0$ for $c \ge c_0$,
closing O5 without any new analytic input.
\end{abstract}

%%=================================================================
\section{The Logical Gap Being Closed}
\label{sec:gap}
%%=================================================================

The three prior O5 files establish:

\begin{enumerate}[label=(\Roman*)]
\item \textbf{Gap:} $\delta_{\mathrm{Airy}} := 1 - \norm{\KAi} \ge F_{\mathrm{TW}}(0) \approx 0.9597$
  (O5-B-1, Tracy--Widom theory).
\item \textbf{Norm convergence:} $\norm{H(c) - I}_{\ell^2 \to \ell^2} \le C_J + \varepsilon(c)$
  with $\varepsilon(c) \to 0$ and $C_J = C\cdot\zeta(3/2)$
  (O5-B-2, Schur/Jaffard + Transfer Lemma).
\item \textbf{Operator norm control:} $\norm{E_c}_{\mathrm{op}} \to 0$
  at rate $(\log\log c / \log c)^{1/2}$
  (O5\_operator\_norm.tex, MV mean square).
\end{enumerate}

What is missing: the deduction of $\inf\sigma(\DAi) > 0$
from (I)--(III).
This is a \emph{spectral stability} conclusion, not a norm estimate,
and it requires an explicit invocation of the appropriate perturbation theorem.
The present file supplies this deduction.

\begin{remark}[Why self-adjointness matters here]
The operator $\DAi = \PAi(\Aarith - \KAi)\PAi$ is self-adjoint
(both $\Aarith$ and $\KAi$ are self-adjoint on $\PAi L^2([0,R])$,
and the projection preserves self-adjointness).
For self-adjoint operators, norm convergence implies spectral convergence
in a precise quantitative sense (Weyl's theorem).
For non-self-adjoint operators this step would require a resolvent argument;
here it is immediate.
\end{remark}

%%=================================================================
\section{Assembling the Perturbation}
\label{sec:assembling}
%%=================================================================

Write $\DAi = T_0 + P_c$ where:
\begin{align}
  T_0 &:= \PAi(I - \KAi)\PAi, \label{eq:T0}\\
  P_c &:= \PAi(\Aarith - I)\PAi = \PAi\cdot E_c\cdot\PAi, \label{eq:Pc}
\end{align}
with $E_c := \Aarith - I$.

\begin{lemma}[Spectrum of $T_0$]\label{lem:T0}
$T_0$ is a non-negative self-adjoint operator on $\PAi L^2([0,R])$,
and its spectral gap satisfies:
\begin{equation}\label{eq:T0gap}
  \inf\sigma(T_0) = 1 - \sup\sigma(\KAi) = 1 - \norm{\KAi} \ge \delta.
\end{equation}
\end{lemma}

\begin{proof}
$\sigma(T_0) = \{1 - \mu : \mu \in \sigma(\KAi|_{\PAi L^2})\}$.
Since $\KAi$ is a positive contraction, $\sigma(\KAi) \subset [0,1]$
and $\sup\sigma(\KAi) = \norm{\KAi} \le 1 - F_{\mathrm{TW}}(0) = 1 - \delta$.
Thus $\inf\sigma(T_0) \ge \delta$. \qed
\end{proof}

\begin{lemma}[Norm of perturbation $P_c$]\label{lem:Pc}
\begin{equation}\label{eq:Pcnorm}
  \norm{P_c}_{\mathrm{op}} \le \eta(c),
\end{equation}
where:
\begin{equation}\label{eq:eta}
  \eta(c) := C_J + \varepsilon(c) + O(c^{-1/3}) \xrightarrow{c\to\infty} C_J.
\end{equation}
The three terms come from:
\begin{enumerate}[label=(\alph*)]
  \item $C_J = C\cdot\zeta(3/2) < F_{\mathrm{TW}}(0)$
        — the Jaffard off-diagonal tail (O4-B-2, transferred via O5 Transfer Lemma);
  \item $\varepsilon(c) \to 0$ — the diagonal convergence rate
        (empirical measure converging to Lebesgue, O5-B-2.1);
  \item $O(c^{-1/3})$ — the basis-identification error
        (Transfer Lemma, Theorem~2.1).
\end{enumerate}
\end{lemma}

\begin{proof}
$P_c = \PAi E_c \PAi$, and $\PAi$ is an orthogonal projection,
so $\norm{P_c}_{\mathrm{op}} \le \norm{E_c}_{\mathrm{op}} = \norm{H(c) - I}_{\ell^2\to\ell^2}$
(in the Airy eigenfunction basis).
The bound $\norm{H(c) - I} \le \eta(c)$ follows from
O5-B-2.3 (Lemma~\ref{lem:coercivity} there) + Transfer Lemma Corollary~3.1.
\qed
\end{proof}

%%=================================================================
\section{The Stability Theorem}
\label{sec:stability}
%%=================================================================

\begin{theorem}[Spectral cluster stability]\label{thm:stability}
Let $c_0$ be large enough that $\eta(c_0) < \delta/2$.
Then for all $c \ge c_0$:
\begin{equation}\label{eq:main}
  \inf\sigma(\DAi) \ge \delta - \eta(c) > 0.
\end{equation}
In particular, $0$ is not in the spectrum of $\DAi$, and
$\DAi$ is boundedly invertible on $\PAi L^2([0,R])$.
\end{theorem}

\begin{proof}
$\DAi = T_0 + P_c$ is self-adjoint (Remark~\ref{sec:gap}).
By Weyl's theorem for self-adjoint operators
(Kato~\cite{Kato1966}, Chapter~IV, Theorem~3.18):
\[
  \inf\sigma(T_0 + P_c) \ge \inf\sigma(T_0) - \norm{P_c}_{\mathrm{op}}.
\]
Substitute Lemma~\ref{lem:T0} and Lemma~\ref{lem:Pc}:
\[
  \inf\sigma(\DAi) \ge \delta - \eta(c).
\]
Since $\eta(c) \to C_J < \delta$ (the numerical condition of O5-B-2,
requiring $C < \delta/\zeta(3/2) \approx 0.595$),
and $c_0$ is chosen so that the transient terms $\varepsilon(c) + O(c^{-1/3}) < \delta - C_J$,
we conclude $\delta - \eta(c) > 0$. \qed
\end{proof}

\begin{remark}[The only numerical condition]
The entire proof reduces to one verifiable inequality:
\begin{equation}\label{eq:numerics}
  C < \frac{F_{\mathrm{TW}}(0)}{\zeta(3/2)} \approx \frac{0.9597}{2.612} \approx 0.367.
\end{equation}
Here $C$ is the off-diagonal oscillation constant from O4-B-2
(the same constant that was already verified numerically in O4).
All other quantities ($F_{\mathrm{TW}}(0)$, $\zeta(3/2)$, $\varepsilon(c)$)
are either classical constants or decay to zero.
\end{remark}

\begin{remark}[Comparison with the non-self-adjoint case]
If $\DAi$ were not self-adjoint, the argument above would require
a resolvent estimate (pseudospectral stability).
Self-adjointness here is not a coincidence: it follows from the fact
that both $\Aarith$ and $\KAi$ are positive-semidefinite sampling/kernel
operators with real, symmetric integral kernels.
The self-adjoint structure is the reason O5 is easier than a
general spectral perturbation problem.
\end{remark}

%%=================================================================
\section{O5: Formal Closure Statement}
\label{sec:closure}
%%=================================================================

\begin{theorem}[O5 closed]\label{thm:O5closed}
Assume:
\begin{enumerate}[label=(\alph*)]
  \item $C < 0.367$ \quad (Numerical condition, pending S4);
  \item MV equidistribution in the Airy window \quad (unconditional, see O5 Remark 4.1).
\end{enumerate}
Then for all $c \ge c_0$ (explicit from PNT + Dusart):
\[
  \inf\sigma\bigl(\PAi(\Aarith - \KAi)\PAi\bigr) > 0,
\]
closing Open Node O5 and proving Conjecture~5.2 of Paper~XXII.
\end{theorem}

\begin{proof}
Combine: Theorem~\ref{thm:stability} (this file)
+ Lemma~O5-B-1 (TW gap)
+ O5 Transfer Lemma (S3 closed)
+ O5-B-2.3 (coercivity of $\Aarith$)
+ O5\_operator\_norm.tex, Theorem~3.4 (MV bound). \qed
\end{proof}

\begin{center}
\renewcommand{\arraystretch}{1.4}
\begin{tabular}{clll}
\toprule
 & \textbf{Node} & \textbf{Status} & \textbf{Blocker?} \\
\midrule
S1 & PNT remainder, explicit $c_0$ & Open (elementary) & No \\
S2 & Prime equidistribution in Airy window & Conditional/MV fallback & No \\
S3 & Basis-change transfer & \textbf{Closed} (Transfer Lemma) & Was yes \\
S4 & Numerical: $C < 0.367$ & Open (computation) & Conditional \\
V1--V4 & Routine analysis & Open (mechanical) & No \\
\textbf{O5} & \textbf{Spectral closure} & \textbf{Closed pending S4} & --- \\
\bottomrule
\end{tabular}
\end{center}

\begin{thebibliography}{99}
\bibitem{Kato1966}
T.~Kato,
\textit{Perturbation Theory for Linear Operators},
Springer, Berlin, 1966.
\bibitem{TracyWidom1994}
C.~Tracy and H.~Widom,
\textit{Level-spacing distributions and the Airy kernel},
Comm.\ Math.\ Phys.\ \textbf{159} (1994), 151--174.
\bibitem{O5Transfer}
S.~Schmalnauer,
\textit{O5\_transfer\_lemma.tex},
\texttt{pswf-programme/core/lemmas/}, 2026.
\bibitem{O5Sampling}
S.~Schmalnauer,
\textit{O5\_B2\_airy\_sampling\_coercivity.tex},
\texttt{pswf-programme/core/open/}, 2026.
\end{thebibliography}

\end{document}
